\documentclass{article}

\title{Footnotes and Endnotes Support}
\author{htex Library}
\date{2026}

\begin{document}

\maketitle

\section{Introduction to Footnotes}

Footnotes are a way to add supplementary information without disrupting the main text.\footnote{This is the first footnote.} They appear at the bottom of the page where they are referenced.

\section{Using Footnotes in Academic Writing}

Academic papers frequently use footnotes for citations and additional context.\footnote{See the bibliography for complete references.} Multiple footnotes can appear on the same page.\footnote{This is another footnote on the same page.}

\subsection{Footnote Examples}

Here are some practical examples:

1. Definitions: A footnote can provide a quick definition.\footnote{A footnote is a note placed at the bottom of a page.}

2. Citations: Alternative to inline citations.\footnote{Smith, J. (2024). "Modern Publishing". Academic Press.}

3. Clarifications: Supplementary details that would clutter the main text.\footnote{The full dataset is available at https://example.com/data}

\section{Academic Paper Example}

The main results show a 45\% improvement\footnote{Measured against baseline from Johnson et al. (2023).} compared to previous work. This was achieved through a novel approach\footnote{Complete algorithm details available in Appendix A.} that leverages recent advances in the field.

Table \ref{table:results} presents the comparative analysis.\footnote{Standard deviation values are shown in parentheses.}

\begin{table}[h]
\centering
\begin{tabular}{|l|r|r|}
\hline
Method & Accuracy & Speed \\
\hline
Baseline & 78.5\% & 2.3s \\
Proposed & 89.2\% & 1.8s \\
\hline
\end{tabular}
\caption{Performance Comparison}
\label{table:results}
\end{table}

\section{Discussion}

The results are statistically significant.\footnote{p < 0.01 using t-test.} Future work should investigate\footnote{We are currently exploring three promising directions.} scaling to larger datasets.

\subsection{Limitations}

Our approach has some limitations.\footnote{GPU memory requirements are approximately 16GB.} These should be considered when applying the method.

\section{Conclusion}

This work demonstrates the effectiveness of the proposed approach.\footnote{Code is available at https://github.com/example/repo} We believe it will inspire future research.\footnote{We welcome collaboration from interested researchers.}

\end{document}
